\chapter{Breadth-First Search and Its Uses}

Breadth-first search is the first algorithm in this book that feels like a
\emph{machine}: you pour a graph and a starting vertex into one end, turn the
crank, and out comes a wave of distances, a tree of shortest routes, and a
linear-time witness for half a dozen structural questions. It is the simplest
graph traversal worth proving things about, and almost every later algorithm --
Dijkstra, the Hopcroft--Karp matching engine, the layered networks of Dinic's
max-flow, even the diameter heuristics on road networks -- is BFS wearing a
costume. Spend real effort here. The invariants you learn to state cleanly for
BFS are the same invariants, dressed up, that carry the heavier chapters.

This chapter builds BFS from the queue up. We start with the picture (a wave
spreading outward), pin down the queue discipline that makes the wave behave,
prove that the wave measures shortest path \emph{in edges}, and then harvest the
applications: the BFS tree, the bipartiteness test, multi-source BFS, and the
0--1 BFS trick that smuggles a tiny bit of weight past the unweighted barrier
using a double-ended queue. Throughout we trace a single small graph by hand so
you can watch the data structures move.

\section{The wave picture}
\label{ch02:wave}

\begin{intuitionbox}
Drop a pebble in still water at vertex $s$. A ring of ripples reaches the
neighbors of $s$ first, then their neighbors, then theirs. The ripple touches a
vertex exactly when the \emph{shortest} chain of edges from $s$ first reaches it,
because the wave moves outward at one edge per tick and never doubles back. BFS
is the bookkeeping that makes a computer reproduce that ripple: it processes
vertices in nondecreasing order of edge-distance from $s$, expanding the closest
unfinished vertex before any farther one.
\end{intuitionbox}

Let $\Gph = (V,E)$ be a graph, directed or undirected, with $n = \abs{V}$
vertices and $m = \abs{E}$ edges, and fix a \emph{source} $s \in V$. Define the
\emph{breadth-first distance} (or unweighted shortest-path distance)
\[
  \dist(s,v) \;=\; \text{the minimum number of edges on any walk } s \to v,
\]
with $\dist(s,v) = \infty$ if no such walk exists. The set of vertices at a fixed
distance $\ell$ is the \emph{level}
\[
  L_\ell \;=\; \set{\, v \in V : \dist(s,v) = \ell \,}.
\]
So $L_0 = \set{s}$, $L_1$ is the neighborhood of $s$, and in general $L_{\ell+1}$
is everything one edge past $L_\ell$ that we have not already seen. The wave
picture is exactly the claim that $V$ partitions into $L_0, L_1, L_2, \dots$ and
that BFS visits the levels in order.

\begin{figure}[H]
\centering
\begin{tikzpicture}[
    every node/.style={circle,draw=nodeblue,thick,fill=blue!5,minimum size=8mm,
      font=\small},
    edge/.style={draw=edgegray,thick},
    level/.style={font=\small\itshape,color=structurepurple}]
  % level 0
  \node (s) at (0,0) {$s$};
  % level 1
  \node (a) at (2.2,1.4) {$a$};
  \node (b) at (2.2,0) {$b$};
  \node (c) at (2.2,-1.4) {$c$};
  % level 2
  \node (d) at (4.6,1.4) {$d$};
  \node (e) at (4.6,-0.3) {$e$};
  \node (f) at (4.6,-1.7) {$f$};
  % level 3
  \node (g) at (7.0,0) {$g$};
  \draw[edge] (s)--(a); \draw[edge] (s)--(b); \draw[edge] (s)--(c);
  \draw[edge] (a)--(d); \draw[edge] (a)--(b);
  \draw[edge] (b)--(e); \draw[edge] (c)--(e); \draw[edge] (c)--(f);
  \draw[edge] (d)--(g); \draw[edge] (e)--(g); \draw[edge] (f)--(g);
  % level brackets
  \node[level,draw=none,fill=none] at (0,-2.4) {$L_0$};
  \node[level,draw=none,fill=none] at (2.2,-2.4) {$L_1$};
  \node[level,draw=none,fill=none] at (4.6,-2.4) {$L_2$};
  \node[level,draw=none,fill=none] at (7.0,-2.4) {$L_3$};
  \begin{scope}[on background layer]
    \fill[blue!4,rounded corners] (-0.7,-2.0) rectangle (0.7,2.0);
    \fill[blue!7,rounded corners] (1.5,-2.0) rectangle (2.9,2.0);
    \fill[blue!10,rounded corners] (3.9,-2.0) rectangle (5.3,2.0);
    \fill[blue!13,rounded corners] (6.3,-2.0) rectangle (7.7,2.0);
  \end{scope}
\end{tikzpicture}
\caption{The running example graph $\Gph_0$ drawn by its BFS levels from $s$.
Each shaded column is a level $L_\ell$; every edge joins vertices in the same
level or in adjacent levels -- never two levels apart. That ``no edge skips a
level'' property is the heart of the correctness proof.}
\label{ch02:fig-levels}
\end{figure}

We will return to $\Gph_0$ from Figure~\ref{ch02:fig-levels} again and again.
For the record, its adjacency lists (undirected, so each edge appears twice) are
\[
\begin{aligned}
  s&:\ a,b,c & a&:\ s,b,d & b&:\ s,a,e & c&:\ s,e,f\\
  d&:\ a,g   & e&:\ b,c,g & f&:\ c,g    & g&:\ d,e,f.
\end{aligned}
\]

\begin{focusbox}
Two facts do all the work in this chapter.
\textbf{(1)} A FIFO queue processes vertices in nondecreasing distance order.
\textbf{(2)} In an unweighted graph every edge connects two vertices whose
distances from $s$ differ by at most $1$. Fact (1) is about the algorithm; fact
(2) is about the graph. Together they force BFS to compute $\dist(s,\cdot)$
exactly.
\end{focusbox}

\section{The queue and the algorithm}
\label{ch02:queue}

The data structure that turns the wave picture into code is the
\emph{first-in-first-out queue}. We push the source, then repeatedly pop a vertex
$u$ from the front, and for each neighbor $v$ of $u$ that we have never seen, we
record its distance as $\dist(u)+1$, mark it seen, and push it to the back. A
vertex is marked the instant it is \emph{discovered} (pushed), not when it is
later processed; marking on discovery is the single most common place beginners
go wrong, and we will dwell on it.

\begin{algorithm}[H]
\DontPrintSemicolon
\SetKwInOut{Input}{input}\SetKwInOut{Output}{output}
\Input{graph $\Gph=(V,E)$ as adjacency lists; source $s\in V$}
\Output{$\dist[v]$ and $\mathrm{parent}[v]$ for every $v$ reachable from $s$}
\ForEach{$v \in V$}{$\dist[v]\gets \infty$; \ $\mathrm{parent}[v]\gets \textsc{nil}$}
$\dist[s]\gets 0$\;
$Q \gets$ empty FIFO queue; \ \textsc{enqueue}$(Q,s)$\;
\While{$Q$ is not empty}{
  $u \gets \textsc{dequeue}(Q)$\;
  \ForEach{$v \in \mathrm{adj}[u]$}{
    \If{$\dist[v] = \infty$}{
      $\dist[v] \gets \dist[u] + 1$\;
      $\mathrm{parent}[v] \gets u$\;
      \textsc{enqueue}$(Q,v)$\;
    }
  }
}
\Return $\dist,\ \mathrm{parent}$\;
\caption{\textsc{Bfs}$(\Gph,s)$ -- breadth-first search from a single source}
\label{ch02:alg-bfs}
\end{algorithm}

Notice that $\dist[v]=\infty$ doubles as the ``not yet discovered'' flag, so we
do not need a separate visited array; the first time we touch $v$ we both stamp
its distance and remove it from the undiscovered pool. The array
$\mathrm{parent}$ records, for each discovered $v$, the vertex from which we first
reached it. Those parent pointers, read backwards, spell out a shortest path, and
collectively they form the \emph{BFS tree} of Section~\ref{ch02:tree}.

\begin{trapbox}
\textbf{Mark on enqueue, not on dequeue.} A tempting but broken variant marks
$v$ as visited only when it is dequeued. Then a single vertex can be pushed many
times before it is ever processed -- once by each neighbor that reaches it while
it sits in the queue. On a dense graph this blows the queue up to $\Th(m)$
entries and, worse, lets a vertex be assigned a larger distance on a later push.
Always set $\dist[v]$ and enqueue in the same breath, guarded by the
$\dist[v]=\infty$ test. One stamp per vertex, forever.
\end{trapbox}

\subsection{Tracing the queue on $\Gph_0$}

Let us run Algorithm~\ref{ch02:alg-bfs} on $\Gph_0$ from $s$, processing each
adjacency list in the listed order. We track the queue contents (front on the
left) and the distance stamps as they are written.

\begin{center}
\renewcommand{\arraystretch}{1.25}
\begin{longtable}{r l l l}
\toprule
step & dequeue & newly discovered (dist) & queue after \\
\midrule
0 & --        & $s\,(0)$                 & $[\,s\,]$ \\
1 & $s$       & $a\,(1),\,b\,(1),\,c\,(1)$ & $[\,a,b,c\,]$ \\
2 & $a$       & $d\,(2)$                  & $[\,b,c,d\,]$ \\
3 & $b$       & $e\,(2)$                  & $[\,c,d,e\,]$ \\
4 & $c$       & $f\,(2)$                  & $[\,d,e,f\,]$ \\
5 & $d$       & $g\,(3)$                  & $[\,e,f,g\,]$ \\
6 & $e$       & --                        & $[\,f,g\,]$ \\
7 & $f$       & --                        & $[\,g\,]$ \\
8 & $g$       & --                        & $[\,\ ]$ \\
\bottomrule
\end{longtable}
\end{center}

\begin{center}
\begin{tikzpicture}[
    >=Stealth,
    queuebox/.style={draw=gray!70, thick, rounded corners=2pt, minimum width=1.05cm, minimum height=0.72cm, font=\small\ttfamily},
    paneltitle/.style={font=\footnotesize\scshape, text=gray!90!black}
]
  % Lane 1: The Queue two-value window
  \node[paneltitle, anchor=west] at (-0.5, 2.2) {Queue contents at intermediate step (Lemma~\ref{ch02:lem-monotone}):};
  
  \node[queuebox, fill=sourceblue!15, draw=sourceblue!80] (q1) at (0.8, 1.3) {$v_1$};
  \node[queuebox, fill=sourceblue!15, draw=sourceblue!80] (q2) at (2.0, 1.3) {$v_2$};
  \node[queuebox, fill=sourceblue!15, draw=sourceblue!80] (q3) at (3.2, 1.3) {$\dots$};
  \node[queuebox, fill=quotientorange!15, draw=quotientorange!80] (q4) at (4.4, 1.3) {$v_i$};
  \node[queuebox, fill=quotientorange!15, draw=quotientorange!80] (q5) at (5.6, 1.3) {$\dots$};
  \node[queuebox, fill=quotientorange!15, draw=quotientorange!80] (q6) at (6.8, 1.3) {$v_r$};

  % Pointers
  \draw[<-, thick, gray!70] (q1.west) -- ++(-0.55, 0) node[left, font=\scriptsize\bfseries] {Front (pop)};
  \draw[->, thick, gray!70] (q6.east) -- ++(0.55, 0) node[right, font=\scriptsize\bfseries] {Back (push)};

  % Braces for distance values
  \node[font=\scriptsize, text=sourceblue!90!black] at (2.0, 0.6) {$\underbrace{\hspace{2.4cm}}_{\dist = k}$};
  \node[font=\scriptsize, text=quotientorange!90!black] at (5.6, 0.6) {$\underbrace{\hspace{2.4cm}}_{\dist = k + 1}$};

  % Lane 2: Outward ripple state
  \node[paneltitle, anchor=west] at (-0.5, -0.1) {Graph partition during traversal:};

  \draw[fill=blue!4, draw=sourceblue!50, rounded corners=3pt] (0.0, -1.0) rectangle (2.1, -0.4);
  \node[font=\scriptsize\bfseries, text=sourceblue!90!black] at (1.05, -0.7) {Settled: $< k$};

  \draw[fill=sourceblue!15, draw=sourceblue!80, rounded corners=3pt] (2.4, -1.0) rectangle (4.9, -0.4);
  \node[font=\scriptsize\bfseries, text=sourceblue!95!black] at (3.65, -0.7) {Active Frontier: $= k$};

  \draw[fill=quotientorange!15, draw=quotientorange!80, rounded corners=3pt] (5.2, -1.0) rectangle (8.0, -0.4);
  \node[font=\scriptsize\bfseries, text=quotientorange!95!black] at (6.6, -0.7) {Next Frontier: $= k + 1$};

  \draw[fill=gray!8, draw=gray!40, rounded corners=3pt] (8.3, -1.0) rectangle (10.3, -0.4);
  \node[font=\scriptsize, text=gray!70!black] at (9.3, -0.7) {Unseen: $\infty$};
\end{tikzpicture}
\end{center}

When $a$ is processed at step 2 it sees neighbors $s$ (already stamped), $b$
(already stamped, it was discovered at step 1), and $d$ (new). So only $d$ is
enqueued. At step 3, $b$ looks at $s,a,e$ and discovers only $e$. The edge $a$--$b$
inside level $L_1$ is examined twice -- once from each endpoint -- and produces no
new discovery either time; those are the \emph{intra-level} edges. The final
distance stamps are
\[
  \dist(s,\cdot):\quad s{=}0,\ a{=}1,\ b{=}1,\ c{=}1,\ d{=}2,\ e{=}2,\ f{=}2,\ g{=}3,
\]
matching the level columns in Figure~\ref{ch02:fig-levels} exactly. The parent
pointers are $\mathrm{parent}[a]=\mathrm{parent}[b]=\mathrm{parent}[c]=s$,
$\mathrm{parent}[d]=a$, $\mathrm{parent}[e]=b$, $\mathrm{parent}[f]=c$, and
$\mathrm{parent}[g]=d$.

\begin{labbox}
Re-run the trace but process each adjacency list in \emph{reverse} order
($s$ expands $c,b,a$). The distances come out identical -- they must, they are a
property of the graph -- but the parent of $g$ changes, and the BFS tree changes
shape. Convince yourself that $\dist[v]$ is canonical while the tree is not.
Then run BFS from $g$ instead of $s$ and check that $\dist(g,s)=3$, confirming
symmetry of distance on undirected graphs.
\end{labbox}

\section{Correctness: BFS computes shortest paths in edges}
\label{ch02:correctness}

We now prove that Algorithm~\ref{ch02:alg-bfs} sets $\dist[v]=\dist(s,v)$ for
every $v$. The argument has two halves that mirror the two facts in the focus
box: a monotonicity lemma about the queue, and the level structure of the graph.

\begin{lemma}[Queue monotonicity]
\label{ch02:lem-monotone}
At every moment during the execution, if the queue holds
$\langle v_1, v_2, \dots, v_r\rangle$ from front to back, then
\[
  \dist[v_1] \le \dist[v_2] \le \cdots \le \dist[v_r]
  \qquad\text{and}\qquad
  \dist[v_r] \le \dist[v_1] + 1 .
\]
That is, distances along the queue are nondecreasing and span a range of at most
one.
\end{lemma}

\begin{proof}
Induct on the number of queue operations. Initially the queue is $\langle s\rangle$
with $\dist[s]=0$; both claims hold vacuously. Suppose they hold before an
operation.

A dequeue removes $v_1$. The remaining queue $\langle v_2,\dots,v_r\rangle$ is
still sorted, and we must recheck the span: $\dist[v_r]\le \dist[v_1]+1 \le
\dist[v_2]+1$, using $\dist[v_1]\le\dist[v_2]$. So the span bound survives.

An enqueue happens only while we are expanding some $u=v_1$ that was just
dequeued, and the new vertex $w$ gets $\dist[w]=\dist[u]+1$. Just before this
expansion the queue was $\langle v_2,\dots,v_r\rangle$ (after $u$ was removed),
with $\dist[v_2]\ge \dist[u]$ (sortedness of the old queue) and
$\dist[v_r]\le \dist[u]+1 = \dist[w]$ (old span bound). Appending $w$ keeps the
sequence sorted because $\dist[w]=\dist[u]+1\ge \dist[v_r]$, and the new span is
$\dist[w]-\dist[v_2] \le (\dist[u]+1) - \dist[u] = 1$. Both invariants are
restored.
\end{proof}

A clean corollary: vertices are dequeued in nondecreasing distance order, and the
distance stamps written over the whole run are nondecreasing in time. The queue
never holds two stamps differing by more than one, which is precisely why BFS
processes one entire level before starting the next.

\begin{lemma}[Adjacent levels]
\label{ch02:lem-adjacent}
For any edge $(u,v)$ of $\Gph$, $\,\abs{\dist(s,u) - \dist(s,v)} \le 1$. For a
directed edge $u\to v$, $\dist(s,v) \le \dist(s,u) + 1$.
\end{lemma}

\begin{proof}
Take any shortest walk $s\to u$ of length $\dist(s,u)$ and append the edge
$u\to v$; this is a walk $s\to v$ of length $\dist(s,u)+1$, so
$\dist(s,v)\le \dist(s,u)+1$. On an undirected graph the edge is also $v\to u$, so
symmetrically $\dist(s,u)\le \dist(s,v)+1$, and the two inequalities give the
absolute-value bound.
\end{proof}

Lemma~\ref{ch02:lem-adjacent} is the formal version of ``no edge skips a level''
in Figure~\ref{ch02:fig-levels}: a single edge can keep you in the same level or
move you one level, never two. With both lemmas in hand the main theorem is a
short induction.

\begin{theorem}[BFS correctness]
\label{ch02:thm-correct}
After Algorithm~\ref{ch02:alg-bfs} halts, $\dist[v]=\dist(s,v)$ for every
$v\in V$, and the value $\infty$ remains exactly for the vertices unreachable
from $s$. Moreover, following parent pointers from any reachable $v$ back to $s$
yields a shortest path.
\end{theorem}

\begin{proof}
First, every distance stamp written by the algorithm is the length of a genuine
$s$--$v$ walk: when $v$ is stamped from $u$ we extend the parent-chain walk of
$u$ by one edge, so by induction the parent chain from $v$ back to $s$ is a real
walk of length $\dist[v]$. Hence $\dist[v]\ge\dist(s,v)$ for every discovered
$v$.

It remains to show that every vertex $v$ with $\dist(s,v)=\ell$ is discovered and
receives $\dist[v]=\ell$. Induct on $\ell$. Base $\ell=0$: only $s$, stamped $0$
before the loop. Inductive step: let $\dist(s,v)=\ell+1$. Some shortest path
reaches $v$ through a vertex $u$ with $\dist(s,u)=\ell$ and edge $u\to v$. By
induction $u$ is discovered with $\dist[u]=\ell$, so $u$ is eventually dequeued.
When $u$ is expanded, either $v$ is still undiscovered -- in which case $v$ is
stamped $\dist[u]+1=\ell+1$ -- or $v$ was already discovered earlier. In the
second case $v$ was stamped from some vertex dequeued no later than $u$. By
Lemma~\ref{ch02:lem-monotone} every vertex dequeued up to and including $u$ has
distance stamp $\le \dist[u]=\ell$, so $v$'s stamp is at most $\ell+1$. Combined
with $\dist[v]\ge\dist(s,v)=\ell+1$ from the first paragraph, the two bounds force
$\dist[v]=\ell+1$. Either way $v$ is discovered with the correct stamp.

Finally, a vertex $v$ is reachable iff $\dist(s,v)$ is finite, and the induction
discovers exactly those. Unreachable vertices are never enqueued, so their stamp
stays $\infty$. The parent chain of a discovered $v$ has lengths decreasing by
exactly one each step down to $s$, so it is a path of length
$\dist[v]=\dist(s,v)$ -- a shortest path.
\end{proof}

\begin{whybox}
The proof never used edge weights because there are none; it used only that a
hop changes distance by one. That is also why BFS \emph{cannot} be salvaged for
weighted graphs by the same argument: with weights, a vertex could be reached
``early but expensively,'' the stamp-once-and-freeze rule would lock in a wrong
value, and Lemma~\ref{ch02:lem-monotone} collapses. Dijkstra's algorithm
(Chapter 7) repairs exactly that crack by replacing the FIFO queue with a
priority queue. BFS is Dijkstra with all weights $1$ and a queue instead of a
heap.
\end{whybox}

\subsection{Complexity}

\begin{theorem}[BFS runs in linear time]
\label{ch02:thm-time}
Algorithm~\ref{ch02:alg-bfs} runs in $\Oh(n+m)$ time and $\Oh(n)$ extra space on
a graph stored as adjacency lists.
\end{theorem}

\begin{proof}
Initialization touches each vertex once: $\Oh(n)$. Each vertex is enqueued at
most once (the $\dist[v]=\infty$ guard fires only on first discovery) and hence
dequeued at most once, so the \textbf{while} loop iterates $\le n$ times with
$\Oh(1)$ queue work apiece. When $u$ is dequeued we scan $\mathrm{adj}[u]$, doing
$\Oh(1)$ per incident edge. Summing over all dequeues, each undirected edge is
scanned from both endpoints (twice) and each directed edge once, for
$\Oh(m)$ total edge work. The grand total is $\Oh(n+m)$. The arrays $\dist$ and
$\mathrm{parent}$ are $\Oh(n)$, and the queue never holds more than $n$ vertices.
\end{proof}

\begin{complexitybox}
\textbf{Single-source BFS.}\quad
Time $\Oh(n+m)$, space $\Oh(n)$. Optimal: any algorithm that even reads a graph
guaranteed to need every edge must spend $\Om(n+m)$. The hidden constant is tiny
-- two array writes and one queue push per discovery -- which is why BFS is the
workhorse for unweighted distance, reachability, and connectivity queries, and
why it is the inner loop of layered max-flow and bidirectional search.
\end{complexitybox}

\begin{historybox}
Breadth-first search was published by Edward F.\ Moore in 1959 as a method for
``finding the shortest path through a maze,'' and independently by C.\ Y.\ Lee in
1961 for routing wires on circuit boards -- the grid version is still called the
\emph{Lee algorithm} in VLSI. Konrad Zuse had described an essentially
breadth-first graph traversal in an unpublished 1945 manuscript. The maze framing
is the right mental model: a maze is a grid graph, every step costs one, and the
ripple from the entrance reaches the exit along a shortest corridor.
\end{historybox}

\section{The BFS tree and shortest-path structure}
\label{ch02:tree}

The parent pointers from Algorithm~\ref{ch02:alg-bfs} define a subgraph
$T=(V_s, E_T)$ on reachable vertices $V_s$, with tree edges
$E_T = \set{(\mathrm{parent}[v],v) : v\ne s,\ v \in V_s}$. Each reachable $v\ne s$ has exactly one
parent and the chain terminates at $s$, so $T$ is a tree rooted at $s$: the
\emph{BFS tree}.

\begin{figure}[H]
\centering
\begin{tikzpicture}[
  every node/.style={circle,draw=nodeblue,thick,fill=blue!5,minimum size=8mm,font=\small},
  tree/.style={draw=targetgreen,very thick},
  nontree/.style={draw=trapred,thick,dashed}]
  \node (s) at (3.5,3) {$s$};
  \node (a) at (1.2,1.6) {$a$};
  \node (b) at (3.5,1.6) {$b$};
  \node (c) at (5.8,1.6) {$c$};
  \node (d) at (0.5,0) {$d$};
  \node (e) at (3.5,0) {$e$};
  \node (f) at (5.8,0) {$f$};
  \node (g) at (2.0,-1.6) {$g$};
  % tree edges (green, solid)
  \draw[tree] (s)--(a); \draw[tree] (s)--(b); \draw[tree] (s)--(c);
  \draw[tree] (a)--(d); \draw[tree] (b)--(e); \draw[tree] (c)--(f);
  \draw[tree] (d)--(g);
  % non-tree edges (red, dashed)
  \draw[nontree] (a) to[bend left=10] (b);
  \draw[nontree] (c) to[bend right=20] (e);
  \draw[nontree] (e) to[bend left=10] (g);
  \draw[nontree] (f) to[bend right=30] (g);
\end{tikzpicture}
\caption{The BFS tree of $\Gph_0$ (green, solid) overlaid with the non-tree
edges (red, dashed). Each non-tree edge is either intra-level ($a$--$b$ joins
two vertices of $L_1$) or joins adjacent levels ($c$--$e$ joins $L_1$ to $L_2$;
$e$--$g$ and $f$--$g$ join $L_2$ to $L_3$). No non-tree edge spans two levels --
the visual restatement of Lemma~\ref{ch02:lem-adjacent}.}
\label{ch02:fig-tree}
\end{figure}

\begin{proposition}[Tree paths are shortest paths]
\label{ch02:prop-treepath}
For every reachable $v$, the unique path from $s$ to $v$ in the BFS tree $T$ has
length $\dist(s,v)$. Consequently $T$ is a \emph{shortest-path tree}: it
simultaneously stores a shortest path from $s$ to every vertex in $\Oh(n)$ space.
\end{proposition}

\begin{proof}
By Theorem~\ref{ch02:thm-correct}, $\dist[v]=\dist[\mathrm{parent}[v]]+1$ for
every reachable $v\ne s$, so the tree depth of $v$ equals $\dist[v]=\dist(s,v)$.
The tree path is therefore a shortest path, and there are at most $n-1$ tree
edges, giving $\Oh(n)$ storage for all $n$ shortest paths at once.
\end{proof}

This is the practical payoff of the whole chapter. To answer ``shortest unweighted
route from $s$ to $v$'' for many targets $v$, run one BFS and then walk parent
pointers; each query costs only the length of its answer. The classification of
non-tree edges in Figure~\ref{ch02:fig-tree} also matters: in an undirected graph
BFS produces only \emph{tree edges} and \emph{cross edges} joining vertices whose
levels differ by $0$ or $1$ -- never a ``back edge'' to a strict ancestor two or
more levels up, and never a ``forward edge'' skipping a level. (DFS, by contrast,
generates back and forward edges; that contrast is the subject of the next
chapter.)

\begin{connectionbox}
The BFS tree is the combinatorial shadow of the graph Laplacian's smoothest
modes. If you order vertices by BFS level and look at the Laplacian $\Lap$ of
$\Gph$, the level structure bounds how fast a function can vary: a vector that
jumps between adjacent levels has bounded Rayleigh quotient, which is one route
into the diameter--eigenvalue inequalities of Chapter 32. We will also reuse BFS
levels directly to build the \emph{layered residual networks} of Dinic's max-flow
(Chapter 12) and the alternating BFS layers of Hopcroft--Karp matching
(Chapter 15). The same wave, three more times.
\end{connectionbox}

\subsection{The eccentricity and diameter view}

The largest finite stamp produced by a BFS from $s$ is the \emph{eccentricity}
$\mathrm{ecc}(s)=\max_v \dist(s,v)$, the distance to the farthest reachable
vertex. The maximum eccentricity over all sources is the \emph{diameter}, the
minimum is the \emph{radius}. On $\Gph_0$, $\mathrm{ecc}(s)=3$ (witnessed by
$g$). Computing exact diameter needs a BFS from every vertex ($\Oh(nm)$ on
sparse graphs), but a two-sweep heuristic -- BFS from any vertex to find a far
vertex $x$, then BFS from $x$ -- gives the \emph{exact} diameter on trees and a
tight estimate on road networks, a trick we revisit in the exercises.

\section{Testing bipartiteness}
\label{ch02:bipartite}

A graph is \emph{bipartite} if its vertices split into two classes so that every
edge crosses between classes -- a $2$-coloring with no monochromatic edge.
Equivalently (we prove this), a graph is bipartite iff it has no odd cycle. BFS
detects bipartiteness for free: color each vertex by the parity of its level.

\begin{intuitionbox}
Walk the wave outward and paint level $L_0$ white, $L_1$ black, $L_2$ white, and
so on -- color $=\dist(s,v) \bmod 2$. Tree edges always cross colors because they
join adjacent levels. The only thing that can ruin the $2$-coloring is a non-tree
edge joining two vertices in the \emph{same} level, because that edge would join
two same-colored vertices. And a same-level edge closes an odd cycle. So the
graph is bipartite iff BFS never finds a same-level edge.
\end{intuitionbox}

\begin{algorithm}[H]
\DontPrintSemicolon
\SetKwInOut{Input}{input}\SetKwInOut{Output}{output}
\Input{undirected graph $\Gph=(V,E)$}
\Output{a $2$-coloring, or report an odd cycle (not bipartite)}
\ForEach{$v\in V$}{$\mathrm{color}[v]\gets \textsc{nil}$}
\ForEach{$s\in V$ \textbf{with} $\mathrm{color}[s]=\textsc{nil}$}{
  $\mathrm{color}[s]\gets 0$; \ $Q\gets\langle s\rangle$\;
  \While{$Q$ not empty}{
    $u\gets\textsc{dequeue}(Q)$\;
    \ForEach{$v\in\mathrm{adj}[u]$}{
      \uIf{$\mathrm{color}[v]=\textsc{nil}$}{
        $\mathrm{color}[v]\gets 1-\mathrm{color}[u]$; \ \textsc{enqueue}$(Q,v)$\;
      }
      \ElseIf{$\mathrm{color}[v]=\mathrm{color}[u]$}{
        \Return \textsc{not bipartite} (odd cycle through $u,v$)\;
      }
    }
  }
}
\Return $\mathrm{color}$\;
\caption{\textsc{BipartiteBfs}$(\Gph)$ -- $2$-color by level parity}
\label{ch02:alg-bipartite}
\end{algorithm}

The outer loop over uncolored $s$ handles disconnected graphs: each component
gets its own BFS, and bipartiteness is a per-component property.

\begin{theorem}[Bipartiteness via BFS]
\label{ch02:thm-bipartite}
$\Gph$ is bipartite iff Algorithm~\ref{ch02:alg-bipartite} colors it without ever
finding a same-color edge, and this runs in $\Oh(n+m)$. Equivalently, $\Gph$ is
bipartite iff it has no odd cycle.
\end{theorem}

\begin{proof}
\emph{If the algorithm succeeds,} the coloring is a valid $2$-coloring by
construction: every edge it accepted joined differently colored endpoints, and it
examined every edge. So $\Gph$ is bipartite, and a bipartite graph has no odd
cycle (a $2$-coloring forces colors to alternate around any cycle, which is
impossible at odd length).

\emph{If the algorithm fails,} it found an edge $(u,v)$ with
$\mathrm{color}[u]=\mathrm{color}[v]$, i.e.\ $\dist(s,u)\equiv\dist(s,v)
\pmod 2$ within the same component. Let $w$ be the lowest common ancestor of $u$
and $v$ in the BFS tree. The closed walk
$w\rightsquigarrow u\to v\rightsquigarrow w$ has length
$(\dist[u]-\dist[w])+1+(\dist[v]-\dist[w]) = \dist[u]+\dist[v]-2\dist[w]+1$,
which is odd because $\dist[u]+\dist[v]$ is even. The two tree paths share only
$w$, so this closed walk contains an odd cycle. An odd cycle blocks any
$2$-coloring, so $\Gph$ is not bipartite.

The two directions together prove the equivalence with ``no odd cycle.'' The cost
is one BFS over each component with an $\Oh(1)$ color check per edge, so
$\Oh(n+m)$.
\end{proof}

\begin{trapbox}
Never eyeball a $2$-coloring from a drawing; run the algorithm. In our running
graph $\Gph_0$, the levels put $a$ and $b$ both in $L_1$, yet the edge $a$--$b$
joins them -- a same-level edge. So $\Gph_0$ is \emph{not} bipartite, and
Algorithm~\ref{ch02:alg-bipartite} certifies the odd cycle $s\to a\to b\to s$ of
length $3$. A single intra-level edge is enough to kill bipartiteness, and a
hand-drawn picture will happily hide it.
\end{trapbox}

\begin{labbox}
Build the $4$-cycle $s\,a\,e\,b\,s$ in isolation (edges $s$--$a$, $a$--$e$,
$e$--$b$, $b$--$s$) and color it from $s$. Levels are $s{=}0$, $a{=}1$, $b{=}1$,
$e{=}2$; every edge crosses a color boundary, so the algorithm succeeds and the
$4$-cycle is bipartite -- as every even cycle is. Then add the chord $a$--$b$ and
watch the same-level edge $a$--$b$ trigger the ``not bipartite'' report, with odd
cycle $s\to a\to b\to s$.
\end{labbox}

\begin{center}
\begin{tikzpicture}[
    >=Stealth,
    vnode/.style={circle, draw=sourceblue, fill=blue!8, thick, minimum size=7.5mm, font=\small\bfseries},
    treeedge/.style={draw=targetgreen!85!black, very thick},
    badedge/.style={draw=trapred!90!black, very thick, dashed},
    colpart/.style={rounded corners=4pt, draw=gray!40, dashed, inner sep=8pt}
]
  % Left partition: Even distance levels (color 0)
  \node[font=\footnotesize\scshape, text=sourceblue!90!black] at (1.5, 3.2) {Class $V_0$ ($\dist \equiv 0 \pmod 2$)};
  \node[vnode] (s) at (1.5, 2.2) {$s$};
  \node[vnode] (d) at (0.6, 0.6) {$d$};
  \node[vnode] (e) at (1.5, -0.6) {$e$};
  \node[vnode] (f) at (2.4, 0.6) {$f$};

  \node[colpart, fill=sourceblue!4, fit=(s) (d) (e) (f), minimum width=3.2cm] (part0) {};

  % Right partition: Odd distance levels (color 1)
  \node[font=\footnotesize\scshape, text=targetgreen!90!black] at (6.5, 3.2) {Class $V_1$ ($\dist \equiv 1 \pmod 2$)};
  \node[vnode] (a) at (5.6, 1.8) {$a$};
  \node[vnode] (b) at (7.4, 1.8) {$b$};
  \node[vnode] (c) at (6.5, 0.6) {$c$};
  \node[vnode] (g) at (6.5, -1.0) {$g$};

  \node[colpart, fill=targetgreen!4, fit=(a) (b) (c) (g), minimum width=3.2cm] (part1) {};

  % Tree edges cross between V_0 and V_1 (valid bipartite cuts)
  \draw[treeedge] (s) -- (a);
  \draw[treeedge] (s) -- (b);
  \draw[treeedge] (s) -- (c);
  \draw[treeedge] (a) -- (d);
  \draw[treeedge] (b) -- (e);
  \draw[treeedge] (c) -- (f);
  \draw[treeedge] (d) -- (g);

  % Intra-level edge inside V_1 creates odd cycle
  \draw[badedge] (a) -- node[midway, above=2pt, font=\scriptsize\bfseries, text=trapred] {monochromatic ($a,b \in L_1$)} (b);

  % Legend below
  \node[anchor=center, font=\scriptsize] at (4.0, -2.1) {
    \textbf{Inspection:}\quad
    \tikz[baseline=-0.5ex]{\draw[treeedge] (0,0) -- (0.5,0);} Bipartite tree edges (cross $V_0 \leftrightarrow V_1$) \qquad
    \tikz[baseline=-0.5ex]{\draw[badedge] (0,0) -- (0.5,0);} Intra-level edge $a$--$b$ certifies odd cycle ($s \to a \to b \to s$)
  };
\end{tikzpicture}
\end{center}

\section{Multi-source BFS}
\label{ch02:multisource}

Sometimes the natural question is not ``how far is $v$ from one source'' but
``how far is $v$ from the \emph{nearest} of several sources.'' Think of fire
stations on a city grid, each cell wanting its distance to the closest station;
or a flood spreading from many wells at once. The trick costs nothing extra:
seed the queue with \emph{all} sources at distance $0$, then run ordinary BFS.

\begin{definition}[Multi-source distance]
\label{ch02:def-multisrc}
Given a set $S\subseteq V$ of sources, define
\[
  \dist(S,v) = \min_{s\in S}\dist(s,v),
\]
the edge-distance from $v$ to the nearest source.
\end{definition}

\begin{algorithm}[H]
\DontPrintSemicolon
\SetKwInOut{Input}{input}\SetKwInOut{Output}{output}
\Input{graph $\Gph=(V,E)$; source set $S\subseteq V$}
\Output{$\dist[v]=\dist(S,v)$ for all $v$}
\ForEach{$v\in V$}{$\dist[v]\gets\infty$}
$Q\gets$ empty FIFO queue\;
\ForEach{$s\in S$}{$\dist[s]\gets 0$; \ \textsc{enqueue}$(Q,s)$}
\While{$Q$ not empty}{
  $u\gets\textsc{dequeue}(Q)$\;
  \ForEach{$v\in\mathrm{adj}[u]$}{
    \If{$\dist[v]=\infty$}{
      $\dist[v]\gets\dist[u]+1$; \ \textsc{enqueue}$(Q,v)$\;
    }
  }
}
\Return $\dist$\;
\caption{\textsc{MultiSourceBfs}$(\Gph,S)$}
\label{ch02:alg-multi}
\end{algorithm}

\begin{whybox}
Why does seeding everything at level $0$ compute the \emph{minimum} over sources?
Build a phantom super-source $s^\star$ joined by a single edge to each $s\in S$,
and run ordinary single-source BFS from $s^\star$. Then $\dist(s^\star,v)=
1+\min_{s\in S}\dist(s,v)$ for $v\notin S$: every shortest path from $s^\star$
must pass through one of its neighbors, namely the nearest source. Stripping off
the leading phantom edge (subtract $1$, and start the real sources at $0$ instead
of $1$) gives exactly Algorithm~\ref{ch02:alg-multi}. Because all of $S$ enters
the queue before any expansion, Lemma~\ref{ch02:lem-monotone} still holds -- the
queue starts as a block of zeros -- so the same monotonicity proof certifies
correctness, and the cost is still $\Oh(n+m)$ regardless of $\abs{S}$.
\end{whybox}

\begin{applicationbox}
\textbf{Grid distance fields.} On an $H\times W$ grid where some cells are sources
(stations, fires, exits), multi-source BFS fills the whole distance field in
$\Oh(HW)$ -- linear in the number of cells, independent of how many sources there
are. This is the standard preprocessing step for influence maps in game AI, for
\emph{distance transforms} in image processing (the city-block distance to the
nearest foreground pixel), and for ``rotting oranges''-style spread simulations.
The phantom-super-source reduction means you never write special-case code: one
BFS, many seeds.
\end{applicationbox}

\section{0--1 BFS with a deque}
\label{ch02:01bfs}

Plain BFS dies the moment edges carry weights, because the FIFO queue assumes
every hop costs exactly one. But there is a sweet spot just past unweighted
graphs where a small modification rescues linear time: graphs whose edge weights
are all $0$ or $1$. Such graphs appear constantly -- a $0$-edge models ``free''
moves (sliding along a conveyor, an already-built road), a $1$-edge a costly move
(building, turning, crossing). For these, replace the queue with a
\emph{double-ended queue} (deque) and the linear-time guarantee survives.

\begin{intuitionbox}
The wave still spreads, but now a $0$-edge lets the wave reach a neighbor
\emph{at no extra distance} -- it belongs to the current ripple, not the next
one. So push $0$-edge neighbors to the \emph{front} of the deque (process them
with the current level) and $1$-edge neighbors to the \emph{back} (process them
with the next level). The deque keeps the same two-value span that ordinary BFS
maintains. It is Dijkstra's algorithm, but the priority queue degenerates into a
deque because there are only two possible key increments.
\end{intuitionbox}

\begin{algorithm}[H]
\DontPrintSemicolon
\SetKwInOut{Input}{input}\SetKwInOut{Output}{output}
\Input{graph $\Gph=(V,E)$ with weights $w(u,v)\in\set{0,1}$; source $s$}
\Output{$\dist[v]=$ shortest path weight $s\to v$}
\ForEach{$v\in V$}{$\dist[v]\gets\infty$}
$\dist[s]\gets 0$; \ $D\gets$ deque containing $s$\;
\While{$D$ not empty}{
  $u\gets\textsc{popFront}(D)$\;
  \ForEach{$v\in\mathrm{adj}[u]$}{
    \If{$\dist[u]+w(u,v) < \dist[v]$}{
      $\dist[v]\gets\dist[u]+w(u,v)$\;
      \uIf{$w(u,v)=0$}{\textsc{pushFront}$(D,v)$}
      \Else{\textsc{pushBack}$(D,v)$}
    }
  }
}
\Return $\dist$\;
\caption{\textsc{ZeroOneBfs}$(\Gph,w,s)$ -- 0--1 BFS with a deque}
\label{ch02:alg-01}
\end{algorithm}

\begin{figure}[H]
\centering
\begin{tikzpicture}[
  every node/.style={circle,draw=nodeblue,thick,fill=blue!5,minimum size=8mm,
    font=\small},
  zero/.style={draw=targetgreen,very thick,-{Stealth[length=2.2mm]}},
  one/.style={draw=trapred,thick,-{Stealth[length=2.2mm]}},
  wlab/.style={font=\scriptsize,draw=none,fill=none,inner sep=1pt}]
  \node (s) at (0,0)   {$s$};
  \node (x) at (2.4,1) {$x$};
  \node (y) at (4.8,0) {$y$};
  \node (t) at (7.2,0) {$t$};
  \draw[one]  (s) -- (x) node[wlab,midway,above left] {$1$};
  \draw[zero] (x) -- (y) node[wlab,midway,above right] {$0$};
  \draw[one]  (y) -- (t) node[wlab,midway,above] {$1$};
  \draw[zero] (s) to[bend right=18] (y);
  \node[wlab] at (2.4,-0.95) {$0$};
\end{tikzpicture}
\caption{The 0--1 BFS example of the lab below. Green arrows are $0$-weight
edges, red arrows are $1$-weight. The free edge $s\to y$ stamps $\dist[y]=0$ by a
front-push before the costly route $s\to x\to y$ can finalize $y$, and $t$
finalizes at total weight $1$.}
\label{ch02:fig-01}
\end{figure}

Two differences from Algorithm~\ref{ch02:alg-bfs} are essential. First, the relax
test is $\dist[u]+w<\dist[v]$, a strict-improvement check, \emph{not} a
discovered-yet flag: a vertex may be reached, then reached again more cheaply via
a chain of $0$-edges, so we must allow re-stamping. Second, because of
re-stamping a vertex can sit in the deque more than once. The strict relax test
makes the duplicates harmless: once $v$ is popped at its final (smallest)
distance, every later relaxation out of a stale copy of $v$ fails the
$<\dist[\cdot]$ guard and changes nothing, so reprocessing it costs only the scan
of its neighbors. Each vertex is finalized once, and the wasted scans are charged
to edges, so the work stays linear.

\begin{theorem}[0--1 BFS correctness and cost]
\label{ch02:thm-01}
On a graph with weights in $\set{0,1}$, Algorithm~\ref{ch02:alg-01} computes
$\dist[v]=\dist_w(s,v)$, the minimum total weight of an $s$--$v$ path, in
$\Oh(n+m)$ time.
\end{theorem}

\begin{proof}
We claim the deque obeys the same monotonicity invariant as the BFS queue:
reading front to back, distances are nondecreasing and span at most $1$. The
invariant holds initially. A $\textsc{popFront}$ removes the minimum-distance
vertex (front), and the span bound is preserved as in
Lemma~\ref{ch02:lem-monotone}. A $\textsc{pushFront}$ adds a vertex $v$ with
$\dist[v]=\dist[u]$ (since $w=0$ and $u$ was the front), so $v$ is a legal new
front. A $\textsc{pushBack}$ adds $v$ with $\dist[v]=\dist[u]+1$, at most one more
than the old back, so it is a legal new back. The span never exceeds $1$.

Consequently vertices are popped in nondecreasing $\dist$ order, exactly the
order Dijkstra would pop them. When a vertex $u$ is popped as the current
minimum, no later relaxation can lower $\dist[u]$ (all remaining deque distances
are $\ge\dist[u]$, and edge weights are nonnegative), so $\dist[u]$ is final and
correct -- the standard Dijkstra finalization argument specialized to the
two-key deque. Hence every $\dist[v]$ converges to $\dist_w(s,v)$.

For the running time, charge each deque insertion to the relaxation that caused
it. Because pops come out in nondecreasing $\dist$ order and weights are in
$\set{0,1}$, when an edge $u\to v$ relaxes successfully the value $\dist[v]$ drops
to $\dist[u]$ or $\dist[u]+1$, and $u$ is then final; a second successful
relaxation of the same edge would require $\dist[u]$ to have dropped, which cannot
happen after $u$ is popped. So each edge contributes $\Oh(1)$ insertions, giving
$\Oh(m)$ insertions and pops overall. Each pop scans the popped vertex's
adjacency list once, for $\Oh(m)$ scan work in total, and initialization is
$\Oh(n)$. Total work $\Oh(n+m)$.
\end{proof}

\begin{labbox}
Take the path $s\xrightarrow{1} x \xrightarrow{0} y \xrightarrow{1} t$ together
with a direct $0$-edge $s\xrightarrow{0} y$, drawn in
Figure~\ref{ch02:fig-01}. Run 0--1 BFS from $s$. You should
watch $y$ get stamped $0$ via the front-push from $s$, \emph{before} the longer
route through $x$ ever finalizes $y$, and then see $t$ finalize at distance $1$
($y\xrightarrow{1}t$). Compare against what plain BFS (treating every edge as
weight $1$) returns for $t$: counting edges, the shortest chain is
$s\to y\to t$ of length $2$, so plain BFS reports $\dist(t)=2$, yet the true
$0$--$1$ weight along that same chain is only $1$ because the $s\to y$ step is
free.
\end{labbox}

\begin{connectionbox}
0--1 BFS is the first rung of a ladder. The next rung, \emph{Dial's algorithm}
(Chapter 8), handles small integer weights $0,1,\dots,C$ with $C+1$ buckets
instead of a two-ended deque; past that lies Dijkstra with a real heap for
arbitrary nonnegative weights (Chapter 7). All three share the finalization
invariant proved above -- pop the current minimum, freeze it -- and differ only
in how cheaply the ``current minimum'' can be extracted. BFS extracts it in
$\Oh(1)$ from a queue, 0--1 BFS in $\Oh(1)$ from a deque, Dijkstra in
$\Oh(\log n)$ from a heap.
\end{connectionbox}

\section{More applications}
\label{ch02:apps}

\paragraph{Connected components and reachability.}
One BFS labels an entire connected component; looping over unvisited vertices and
launching a BFS from each labels every component in $\Oh(n+m)$. On a directed
graph a single BFS from $s$ computes the set of vertices reachable \emph{from}
$s$; running BFS on the reverse graph computes the set that can \emph{reach} $s$.
Their intersection is used in Kosaraju's strongly-connected-components routine of
Chapter 5.

\paragraph{Shortest cycle (girth) through unweighted BFS.}
To find the shortest cycle in an undirected unweighted graph, run a BFS from each
vertex $r$; if while expanding you ever find a non-tree edge $(u,v)$ with both
endpoints already discovered in the \emph{current} BFS, it closes a cycle of
length $\dist[u]+\dist[v]+1$. The minimum over all roots and all such edges is the
girth, in $\Oh(nm)$.

\paragraph{Word ladders and implicit graphs.}
BFS does not need the graph in memory. In a word-ladder puzzle (\texttt{cat} to
\texttt{dog}, changing one letter at a time through valid words), vertices are
words and edges connect words differing in one letter. Generate neighbors on the
fly inside the inner loop; BFS finds the fewest-step chain. The same
\emph{implicit-graph} pattern solves sliding-tile puzzles, Rubik's-cube searches
(with bidirectional BFS), and state-space planning.

\paragraph{Bidirectional BFS.}
For a single $s$--$t$ query on a huge graph, grow two BFS waves, one from $s$ and
one from $t$, alternating and stopping when they meet. If the graph has branching
factor roughly $b$ and the answer is at distance $\ell$, one-directional BFS
touches $\Th(b^\ell)$ vertices while the two half-depth waves touch only
$\Th(b^{\ell/2})$ each -- a square-root savings that turns hopeless searches
tractable. The subtlety is the meeting condition: stop only once the combined
frontier radii can no longer improve, then reconstruct the path by splicing the
two parent chains at the cheapest meeting vertex.

\begin{applicationbox}
\textbf{Maze and grid routing, restated.} Moore's and Lee's original use cases
are still live. A robot on an occupancy grid runs multi-source BFS from all
obstacle cells to get a clearance map, then single-source BFS on the free cells
for a shortest collision-free route. The $0$--$1$ variant models terrain where
some moves are free (downhill, paved) and others cost one (rough), all without
leaving linear time. When weights become genuinely varied, you graduate to
Dijkstra -- but a surprising amount of real routing lives in the $\set{0,1}$ and
small-integer regime where BFS-family algorithms win on constant factors.
\end{applicationbox}

\section{Summary}
\label{ch02:summary}

BFS expands vertices in nondecreasing edge-distance from the source using a FIFO
queue, runs in $\Oh(n+m)$, and produces a shortest-path tree as a side effect.
The two load-bearing facts are queue monotonicity
(Lemma~\ref{ch02:lem-monotone}) and the adjacent-level property of unweighted
graphs (Lemma~\ref{ch02:lem-adjacent}); together they yield the correctness
theorem and every application. Color by level parity to test bipartiteness; seed
the queue with many vertices for multi-source distance fields; swap the queue for
a deque to handle $\set{0,1}$ weights. Everywhere in this book that a layered
structure appears -- Dinic, Hopcroft--Karp, the Laplacian's smooth modes -- the
BFS wave is underneath it. Master the invariants here once and you read those
chapters at a glance.

\begin{complexitybox}
\textbf{The chapter on one line.}\quad
Single/multi-source BFS: $\Oh(n+m)$ time, $\Oh(n)$ space, exact unweighted
distance and a shortest-path tree. Bipartiteness: $\Oh(n+m)$. 0--1 BFS:
$\Oh(n+m)$ for weights in $\set{0,1}$ via a deque. Bidirectional BFS:
$\Th(b^{\ell/2})$ for a single far query.
\end{complexitybox}

\section{Exercises}
\label{ch02:exercises}

\begin{exercise}[Implement and instrument]
Implement Algorithm~\ref{ch02:alg-bfs} and, additionally, return the level sizes
$\abs{L_0},\abs{L_1},\dots$. Run it on $\Gph_0$ and confirm the sizes
$1,3,3,1$. Then add a counter for how many edges are examined; verify it equals
$2m$ on an undirected graph and explain the factor of two.
\end{exercise}

\begin{exercise}[Distances are canonical, trees are not]
Prove that for a fixed source $s$, the array $\dist[\cdot]$ produced by BFS is
independent of the order in which adjacency lists are scanned, but the
$\mathrm{parent}[\cdot]$ array (hence the BFS tree) need not be. Give two
distinct BFS trees of $\Gph_0$ from $s$.
\end{exercise}

\begin{exercise}[Break the mark-on-dequeue bug]
Modify Algorithm~\ref{ch02:alg-bfs} to mark vertices visited only when dequeued
(the trap in Section~\ref{ch02:queue}). Construct a family of graphs on $n$
vertices for which this buggy version enqueues $\Th(n^2)$ times, and exhibit a
vertex that holds an incorrect (too large) distance at some intermediate moment.
Then prove the final distances are still correct despite the blowup, explaining
why this is a \emph{performance} bug, not always a correctness bug.
\end{exercise}

\begin{exercise}[Two-sweep diameter on trees]
Prove that on a \emph{tree}, the following computes the exact diameter: BFS from
any vertex to find a farthest vertex $x$; BFS from $x$ to find a farthest vertex
$y$; return $\dist(x,y)$. Show by counterexample that the same two-sweep method
can underestimate the diameter on a general graph with a cycle.
\end{exercise}

\begin{exercise}[Odd-cycle certificate]
Extend Algorithm~\ref{ch02:alg-bipartite} so that when it reports ``not
bipartite,'' it returns an explicit odd cycle as a vertex list. Prove your cycle
has odd length and bound its length by $2\,\mathrm{ecc}(s)+1$.
\end{exercise}

\begin{exercise}[Multi-source distance transform]
You are given an $H\times W$ grid with obstacle cells. Using one multi-source BFS,
compute for every free cell its distance to the nearest obstacle (the city-block
distance transform). State the running time and prove correctness via the
phantom super-source reduction of Section~\ref{ch02:multisource}.
\end{exercise}

\begin{exercise}[0--1 BFS on a grid with portals]
On a grid, moving to an orthogonal neighbor costs $1$, but stepping onto a portal
teleports to its twin at cost $0$. Model this as a $\set{0,1}$-weighted graph and
give the shortest-travel-cost algorithm. Prove that your deque-based solution runs
in $\Oh(HW)$ and argue why a plain FIFO queue would give wrong answers.
\end{exercise}

\begin{exercise}[Lower bound on insertions]
Prove that 0--1 BFS performs $\Oh(m)$ deque insertions even with re-stamping, by
charging each insertion to an edge relaxation that strictly decreased some
$\dist[v]$. Conclude the $\Oh(n+m)$ bound rigorously, addressing the stale-entry
skips.
\end{exercise}

\begin{exercise}[Bidirectional correctness]
Prove that bidirectional BFS returns a shortest $s$--$t$ path, being careful about
the stopping rule: show that halting at the \emph{first} vertex discovered by both
frontiers can overshoot, and give the corrected rule (continue until the sum of
the two frontier radii can no longer improve) with a proof.
\end{exercise}

\begin{exercise}[BFS and the Laplacian, a first taste]
Let $\Gph$ be connected with BFS levels $L_0,\dots,L_k$ from $s$. Define
$x\in\R^V$ by $x_v=\dist(s,v)$. Writing $\Lap=\mathbf{D}-\Adjm$, show that
$(\Lap x)_v$ equals the number of neighbors of $v$ in lower levels minus the
number in higher levels (same-level neighbors cancel). Compute the Rayleigh
quotient $x^\top \Lap x / x^\top x$ for $\Gph_0$, and relate its magnitude to the
eccentricity, foreshadowing Chapter 32.
\end{exercise}

\begin{exercise}[Word ladder]
Given a dictionary, build the implicit one-letter-change graph and use BFS to find
the shortest transformation from a start word to a target word. Analyze the cost
in terms of word length $L$, alphabet size $\sigma$, and dictionary size $N$, and
explain why generating neighbors on the fly beats materializing all $\Th(N^2)$
potential edges.
\end{exercise}
